\chapter{Análise dos resultados}
\label{cap:resultados}

Este capítulo distingue evidências de validação funcional, calibração dos instrumentos e resultados comparativos. A primeira etapa demonstra os comportamentos observados do protótipo e a rastreabilidade da coleta, mas não substitui as repetições previstas no método. A resposta às hipóteses de desempenho e resiliência depende de uma campanha experimental completa, com parâmetros congelados e análise entre execuções independentes. Nesta versão, as evidências disponíveis são identificadas pelo seu escopo, sem extrapolar o comportamento local para REST ou Kafka em geral.

\section{Ferramentas de carga e piloto}

Foram examinadas as documentações oficiais do Apache JMeter e do Grafana k6. O Quadro~\ref{qua:ferramentas} registra a comparação funcional utilizada na escolha do instrumento. Não foi executado um benchmark entre os geradores e, portanto, a escolha não afirma que k6 tenha melhor desempenho que JMeter.

\begin{quadro}[htbp]
\centering
\caption{Comparação dos instrumentos de geração de carga}
\label{qua:ferramentas}
\small
\begin{tabular}{|p{0.19\textwidth}|p{0.35\textwidth}|p{0.35\textwidth}|}
\hline
\textbf{Critério} & \textbf{Apache JMeter} & \textbf{Grafana k6} \\ \hline
Plano & Plano de teste JMX, interface gráfica e propriedades parametrizáveis. & Funções e cenários definidos em JavaScript. \\ \hline
Execução & Modo de linha de comando recomendado para carga; listeners gráficos aumentam o consumo. & Linha de comando e execução em contêiner. \\ \hline
Chegada aberta & Open Model Thread Group, classificado na documentação como experimental. & Executores de taxa constante e variável, independentes da duração das respostas. \\ \hline
Dados de saída & JTL/CSV; recomenda-se limitar os dados coletados. & JSON por linha, comprimível, com tags de correlação. \\ \hline
\end{tabular}
\fonte{Elaborado pelo autor com base nas documentações oficiais das ferramentas, consultadas em 6 out. 2026.}
\end{quadro}

As orientações de execução em linha de comando e redução de listeners constam das boas práticas do JMeter\footnote{\url{https://jmeter.apache.org/usermanual/best-practices.html}}. A existência de um modelo aberto também no JMeter\footnote{\url{https://jmeter.apache.org/usermanual/component_reference.html}, seção Open Model Thread Group.} evita atribuir essa capacidade exclusivamente ao k6. No k6, o executor de chegada constante inicia iterações segundo uma taxa programada, desde que existam usuários virtuais disponíveis\footnote{\url{https://grafana.com/docs/k6/latest/using-k6/scenarios/executors/constant-arrival-rate/}}. A saída JSON permite preservar amostras e metadados por evento\footnote{\url{https://grafana.com/docs/k6/latest/results-output/real-time/json/}}. Essas fontes são documentação técnica dos instrumentos; não substituem os artigos científicos discutidos na revisão bibliográfica.

O k6 foi escolhido pela integração com o protocolo já implementado: scripts versionados permitem gerar os UUIDs e os conteúdos sintéticos, separar aquecimento e medição, registrar a fase de cada amostra e enviar o identificador da execução no cabeçalho HTTP. A imagem utilizada é a versão 2.3.0, fixada por digest. O analisador Java cruza os dados do gerador com os marcos de transação e os registros exportados do PostgreSQL. A exportação por evento e a instrumentação também consomem recursos, razão pela qual precisam permanecer constantes nas comparações e ser reconhecidas como limitação.

\subsection{Validação inicial dos instrumentos}

Em 5 de outubro de 2026, foram realizadas duas execuções curtas válidas de conferência, à taxa de 2 eventos por segundo durante 10 segundos. A variante síncrona registrou 21 eventos e a assíncrona, 20. Em cada execução, todos os eventos observados foram associados a respostas esperadas, marcos de confirmação e registros com conteúdo e hash correspondentes. A diferença de uma iteração nos limites da execução reforça a necessidade de informar a carga efetiva, além da programada. Esses ensaios não usaram o mesmo conjunto entre as variantes e não estimam a carga de referência.

Cinco tentativas anteriores foram preservadas como inválidas por falha de conciliação dos logs do ambiente Docker, embora as linhas tivessem sido encontradas no banco. O problema foi investigado antes de aceitar novas coletas. Não se interpretou ausência de marco como perda de registro. O histórico de tentativas e seus UUIDs está preservado no relatório técnico do piloto e nas evidências locais.

O intervalo nativo de resposta HTTP do k6 é separado do intervalo instrumentado entre envio e retorno. A conclusão da persistência é observada pelo marco \texttt{commit\_completed}, emitido após o retorno da transação. O campo \texttt{persisted\_at}, preenchido no INSERT, não é utilizado como substituto desse marco. Uma resposta 202 confirma publicação no Kafka e não gravação no PostgreSQL.

\section{Modelo comum de registros: objetivo a}
\label{sec:resultado-modelo}
Em 6 de outubro de 2026, o teste de equivalência enviou um conjunto determinístico de 50 registros a cada variante, em execuções separadas. Os UUIDs, campos, instantes de ocorrência e payloads foram iguais; somente o identificador de execução, transportado fora do conteúdo canônico, diferiu. Ambas as exportações continham 50 linhas. A comparação independente confirmou igualdade dos identificadores, campos, payloads e SHA-256, desconsiderando apenas \texttt{persisted\_at}, pois as inserções ocorreram em momentos distintos.

O reenvio do primeiro registro, com o mesmo conteúdo, produziu um marco de persistência \texttt{duplicate} em cada variante e não aumentou a contagem de linhas. Em seguida, o payload desse UUID foi alterado de zero para nove. A API síncrona respondeu 409. A assíncrona respondeu 202 por confirmar a publicação; o consumidor identificou \texttt{event\_id\_content\_conflict} e confirmou o encaminhamento da mensagem conflitante à DLQ. O conteúdo original no banco permaneceu inalterado.

Esse resultado evidencia equivalência funcional para o conjunto examinado e verifica os caminhos de idempotência e conflito. Não demonstra equivalência para todos os possíveis payloads ou condições de concorrência e não constitui comparação de desempenho. As execuções válidas foram identificadas pelos prefixos \texttt{e71d7147} e \texttt{6cdf03fd}; os UUIDs completos, as exportações e seus checksums constam do catálogo técnico de evidências.


\section{Variante síncrona: objetivo b}
A implementação síncrona recebe o evento pela API REST e conclui a transação antes de devolver a resposta de sucesso. A validação funcional pode conferir o processamento normal, a repetição idempotente e a indisponibilidade do banco. A caracterização quantitativa sob carga baixa, crescente e durante falhas exige as repetições da campanha definitiva; não se infere essa caracterização a partir de ensaios curtos.

\section{Variante assíncrona: objetivo c}
O aceite e a persistência são etapas distintas. Com Kafka disponível, a API pode confirmar a publicação enquanto o consumidor está indisponível. A interrupção do consumidor difere da interrupção do PostgreSQL: no segundo caso, o consumidor ativo pode esgotar suas tentativas e enviar eventos à DLQ. A retomada do banco não reprocessa automaticamente essa fila. A análise deverá informar tanto o tempo até o aceite quanto o tempo até a confirmação de gravação e o destino dos eventos que não foram persistidos.

\section{Execução dos cenários: objetivo d}
O executor experimental cria um projeto Docker Compose separado, sem utilizar as portas ou os volumes da aplicação interativa. Cada ensaio registra sua configuração, seus UUIDs e sua classificação. Os testes curtos de validação não entram na contagem de repetições definitivas. Os cenários planejados e os critérios prévios permanecem descritos no método; o catálogo de execuções distinguirá o que foi efetivamente realizado, as tentativas inválidas e seus motivos.

Na validação com PostgreSQL interrompido, foi programada carga de 5 eventos por segundo, com 5 segundos de aquecimento, intervalo de até 12 segundos para concluir o aquecimento, 18 segundos de medição e 5 segundos de interrupção solicitada ao controlador. O banco foi interrompido aproximadamente no primeiro terço da medição. Os instantes dos comandos de parada, da confirmação de parada e do comando de retomada foram registrados separadamente, pois a duração operacional pode diferir da pausa programada. Após a carga, foi aplicada uma observação de 5 segundos, além do tempo necessário à exportação.

A execução síncrona observou 91 eventos na fase medida: 39 respostas de sucesso e 52 respostas 503; o banco continha 39 desses registros ao final. A assíncrona observou 90 eventos: todos receberam 202, mas somente 89 estavam no banco, e um foi encontrado na DLQ por esgotamento das tentativas de persistência. Não havia aceite não conciliado nessas coletas. A mensagem na DLQ não desapareceu, mas também não equivale a um registro salvo. A diferença de uma iteração na fronteira das execuções é explicitada e impede tratar os totais como volumes exatamente iguais.

O comportamento observado é compatível com a distinção entre confirmação de publicação e conclusão da persistência. Ele impede afirmar que o aceite assíncrono garante que todo registro será salvo automaticamente depois que o banco voltar. Trata-se, contudo, de uma única validação curta por variante, não de uma estimativa de taxa de erro, capacidade ou recuperação generalizável. Não se confirmam ou refutam hipóteses comparativas de desempenho com esse ensaio.

Os perfis de rajada também foram exercitados como validação: a taxa passou de 5 para 10 eventos por segundo e retornou a 5, em uma janela de 18 segundos. Foram observados 120 registros na síncrona e 119 na assíncrona; todos esses registros foram conciliados com o banco, sem divergência de conteúdo ou hash. Os totais diferem nas fronteiras do agendamento e não representam uma comparação estatística de capacidade.

Na primeira validação com somente o consumidor interrompido, 91 eventos receberam 202, mas a observação inicial ainda encontrou 54 pendentes retidos e 37 salvos. Não se declarou perda. Um novo ensaio do mesmo perfil, com observação posterior de 40 segundos, conciliou todos os 91 eventos no banco, sem DLQ ou pendência final. Nesse segundo ensaio, 51 registros aceitos estavam sem confirmação de persistência no instante de conclusão do comando de retomada. O primeiro commit posterior foi observado aproximadamente 9,83 segundos após esse comando, e a persistência de todo esse conjunto pendente foi confirmada aproximadamente 10,44 segundos depois. Esses tempos não medem estabilidade sustentada de throughput e não são médias entre repetições. O aumento da janela de observação pertence à validação prévia dos instrumentos, não a uma alteração posterior de protocolo congelado.

A primeira tentativa de exportação do Kafka utilizou um consumidor de grupo cujo prazo de leitura terminou durante a atribuição inicial, produzindo exportação vazia apesar da confirmação da DLQ nos logs. A coleta foi rejeitada por inconsistência do instrumento e preservada. A exportação passou a atribuir diretamente a partição zero; um novo teste confirmou o conteúdo. Uma execução síncrona anterior também foi interrompida por alteração do script durante sua execução. Essas tentativas não foram incluídas como evidências válidas nem apagadas do histórico técnico.


\section{Métricas comparativas: objetivo e}
Ainda não existem as dez repetições por combinação de variante e cenário previstas no método. Por esse motivo, não são apresentadas tabelas de desempenho como se os cenários definitivos já tivessem sido executados. O analisador passou a calcular distribuições por execução, throughput concluído dentro da janela de medição, categorias de resposta, destinos dos eventos e recursos por contêiner. Salvar um evento durante a drenagem não aumenta o throughput da janela anterior.

A campanha deverá agregar resultados por repetição e apresentar dispersão, carga efetiva e quantidade de observações. O número de eventos não representa o número de repetições independentes. Recuperação deve separar o primeiro commit após a retomada, a conclusão do backlog e a retomada de throughput estável. Caso existam eventos na DLQ, não cabe declarar recuperação completa da persistência somente porque o lag do grupo chegou a zero.

\section{Vantagens, desvantagens e hipóteses: objetivo f}
As hipóteses comparativas permanecem inconclusivas na ausência da campanha completa. O sucesso dos testes funcionais demonstra apenas os comportamentos examinados no protótipo. A menor quantidade de componentes no caminho síncrono não mede, por si só, custo de operação; tampouco uma resposta 202 mais rápida demonstra conclusão mais rápida do trabalho. As recomendações entre as abordagens serão delimitadas pelas condições efetivamente medidas, inclusive se os resultados contrariarem as expectativas iniciais.

\section{Limitações das evidências disponíveis}
O protótipo e o gerador compartilham um computador Windows com Docker em WSL. Há um único broker, uma partição e fator de replicação um; o aceite do broker não constitui teste de tolerância à perda da máquina. Os dados são sintéticos e possuem payload reduzido. A telemetria adiciona custo e as leituras de CPU e memória são amostras do contêiner, não consumo global do computador nem custo financeiro de manutenção.

Os intervalos entre processos dependem de relógios de parede. Intervalos negativos são detectáveis, mas a ausência deles não elimina a possibilidade de desvio. Além disso, componentes concorrentes e o estado da virtualização podem influenciar os ensaios. Essas restrições impedem generalizar o resultado local para todas as aplicações REST, todas as configurações Kafka ou ambientes de produção.








